\documentclass[13pt, a4paper]{extarticle} % Hỗ trợ cỡ chữ 13pt chuẩn BM18

% ----- CÀI ĐẶT FONT VÀ NGÔN NGỮ (DÙNG PDFLATEX) -----
\usepackage[utf8]{inputenc} % Hỗ trợ gõ tiếng Việt
\usepackage[T5]{fontenc}    % Bảng mã tiếng Việt
\usepackage{mathptmx}       % Font Times mổ phỏng cho pdfLaTeX
\usepackage[vietnamese]{babel}

% ----- CÀI ĐẶT KHỔ GIẤY VÀ LỀ -----
\usepackage{geometry}
\geometry{
    a4paper,
    left=3cm, right=2cm, top=2.5cm, bottom=2.5cm % Lề tiêu chuẩn BM18
}

% ----- CÀI ĐẶT ĐÁNH SỐ TRANG GIỮA BÊN TRÊN -----
\usepackage{fancyhdr}
\pagestyle{fancy}
\fancyhf{} % Xóa mọi định dạng header/footer cũ
\fancyhead[C]{\thepage} % Đánh số trang ở chính giữa phía trên
\renewcommand{\headrulewidth}{0pt} % Xóa đường kẻ ngang ở header

% ----- CÁC GÓI TOÁN HỌC VÀ HÌNH ẢNH DÀNH CHO VRPTW -----
\usepackage{amsmath, amssymb}
\usepackage{graphicx}
\usepackage{caption}
\usepackage{indentfirst} % Thụt lề đoạn đầu tiên

% ----- CÀI ĐẶT TRÍCH DẪN (AUTHOR-YEAR) -----
\usepackage[authoryear, round]{natbib}
\renewcommand{\bibname}{Tài liệu tham khảo}

\begin{document}

% ================= TRANG BÌA =================
\begin{titlepage}
    \thispagestyle{empty} % TẮT ĐÁNH SỐ TRANG CHO TRANG BÌA
    \begin{center}
        TỔNG LIÊN ĐOÀN LAO ĐỘNG VIỆT NAM \\
        TRƯỜNG ĐẠI HỌC TÔN ĐỨC THẮNG \\
        --------------------------------- \\
        \vspace{2.5cm}
            ĐÁNH THỬ, COI NHƯ NHÁP, KHÔNG DÙNG CÁI NÀY ĐƯA CÔ COI, NĂN NỈ NĂN NỈ NĂN NỈ
            \vspace{1.5cm}
        
        Công trình Nghiên cứu khoa học sinh viên năm học 2025 - 2026 \\
        \vspace{1.5cm}
        
        TÊN ĐỀ TÀI \\
        \vspace{0.5cm}
        \textbf{\Large TỐI ƯU HÓA BÀI TOÁN VRPTW SỬ DỤNG MẠNG NƠ-RON HỌC TĂNG CƯỜNG SÂU KẾT HỢP TÌM KIẾM LÂN CẬN LỚN THÍCH ỨNG (DDQN-ALNS)} \\
        \vspace{1.5cm}
        
        ĐƠN VỊ: KHOA CÔNG NGHỆ THÔNG TIN
        \vspace{2.5cm}
    \end{center}
    
    \begin{flushleft}
        GIẢNG VIÊN HƯỚNG DẪN: \\
        TS. Hồ Thị Linh \\
        \vspace{0.8cm}
        SINH VIÊN THỰC HIỆN: \\
        1. Huỳnh Nhật Huy \\
        2. Nguyễn Nhật Huy \\
        3. .......... Trân

    \end{flushleft}
    
    \vfill
    \begin{center}
        TP. Hồ Chí Minh, tháng 4 năm 2026
    \end{center}
\end{titlepage}

\newpage
\setcounter{page}{1}

%========================================================== INTRODUCTION
\section*{GIỚI THIỆU VÀ ĐẶT VẤN ĐỀ}
Trong bối cảnh logistics và thương mại điện tử hiện nay, việc tối ưu hóa lộ trình giao hàng không chỉ giúp tiết kiệm hàng tỷ đồng chi phí nhiên liệu mà còn đáp ứng yêu cầu khắt khe về thời gian của khách hàng. Bài toán VRPTW (Vehicle Routing Problem with Time Windows) ra đời để giải quyết bài toán này. Tuy nhiên, việc sử dụng các heuristic tĩnh truyền thống để giải quyết đồ thị lớn thường dẫn đến hiện tượng thuật toán bị kẹt ở các cực trị địa phương, không thể tiếp tục tối ưu hóa. Điều này đặt ra nhu cầu cấp thiết về một bộ điều khiển thông minh, có khả năng tự động điều chỉnh chiến lược tìm kiếm theo thời gian thực.
\newpage

%========================================================== 

\section{Tổng quan tài liệu}
Các thuật toán như Quy hoạch tuyến tính của Dantzig và Ramser (1959) đã đặt nền móng đầu tiên cho VRP. Gần đây, phương pháp Adaptive Large Neighborhood Search (ALNS) trở thành tiêu chuẩn nhờ khả năng linh hoạt trong việc phá hủy và tái thiết đồ thị. Dù vậy, ALNS thuần túy bị hạn chế bởi việc chọn toán tử dựa trên xác suất mù quáng. Việc ứng dụng Học tăng cường sâu (Deep Reinforcement Learning) để điều khiển metaheuristics \citep{Author2023} đang mở ra hướng đi mới đầy hứa hẹn, cho phép hệ thống học được các "siêu chiến lược" (meta-policy) từ dữ liệu mô phỏng.

\section{Mục tiêu - Phương pháp}

\subsection{Mục tiêu nghiên cứu}
Mục tiêu chính là xây dựng bộ giải \textit{PlateauHybridSolver} sử dụng mạng DDQN nhằm khắc phục tình trạng chững lại của ALNS, đồng thời ngăn chặn hiện tượng thuật toán đánh đổi việc tăng số lượng xe để giảm quãng đường di chuyển.

\subsection{Phương pháp nghiên cứu (Methodology)}

Kiến trúc DDQN-ALNS được mô hình hóa dưới dạng một Quá trình Quyết định Markov (MDP) với ba thành phần cốt lõi: Không gian trạng thái, Không gian hành động, và Hàm phần thưởng. Đặc điểm nổi bật của hệ thống là cơ chế nhận diện điểm chững: DDQN chỉ can thiệp khi ALNS không cải thiện kết quả sau 60 vòng lặp liên tiếp.

\subsubsection{Không gian Trạng thái (State Space)}
Để tác nhân DDQN có thể quan sát cấu trúc và tình trạng tìm kiếm hiện tại mà không bị phụ thuộc vào quy mô đồ thị, một vector trạng thái gồm 10 chiều được trích xuất và chuẩn hóa về khoảng [0, 1]. Phương trình trạng thái $S_t$ được định nghĩa như sau:

\[
S_t = \left[ \text{patience}, \text{gap}, \text{temp\_ratio}, \text{imp\_rate}, \text{nv\_ratio}, \text{route\_balance}, \text{load\_balance}, \text{tw\_tightness}, \text{avg\_slack}, \text{progress} \right]
\]

Các biến số phản ánh độ bế tắc hiện tại (patience), sự chênh lệch so với nghiệm tốt nhất (gap), trạng thái cân bằng tải (load\_balance), và mức độ nghiêm ngặt của các ràng buộc thời gian (tw\_tightness).

\subsubsection{Không gian Hành động (Action Space)}
Thay vì lựa chọn trực tiếp các toán tử cơ sở, DDQN được thiết kế để điều hướng không gian tìm kiếm bằng cách đưa ra 1 trong 4 chiến lược tổng quát (Modes):
\begin{itemize}
    \item \textbf{Mode 0 (Default):} Tiếp tục tỷ lệ xác suất mặc định của ALNS.
    \item \textbf{Mode 1 (Intensify):} Khai thác sâu vùng không gian hiện tại, ưu tiên các toán tử loại bỏ khắt khe như \textit{Worst} hoặc \textit{Shaw removal}.
    \item \textbf{Mode 2 (Diversify):} Phá vỡ cấu trúc nghiệm hiện tại bằng cách tăng cường hệ số phá hủy (destroy\_scale = 1.35) nhằm nhảy khỏi cực trị địa phương.
    \item \textbf{Mode 3 (TW Rescue):} Tập trung xử lý các điểm giao hàng đang vi phạm nghiêm trọng cửa sổ thời gian.
\end{itemize}

\subsubsection{Hàm Phần thưởng nhận thức sức chứa (Capacity-aware Reward)}
Trong tối ưu hóa VRPTW, số lượng xe (NV) đóng vai trò ưu tiên cao hơn tổng quãng đường (TD). Các phương pháp cũ thường gặp "hội chứng trôi dạt", tự ý phân bổ thêm xe để làm giảm TD. Nghiên cứu này đề xuất một hàm phần thưởng $R_t$ áp đặt mức phạt nghiêm khắc (NV Penalty = -15.0) để triệt tiêu hành vi này. Giả sử $\Delta NV = NV_{after} - NV_{before}$ và $\Delta TD = TD_{before} - TD_{after}$ (suy ra $\Delta TD > 0$ là quãng đường đã được rút ngắn), hàm phần thưởng được biểu diễn như sau:

\[
R_t = -0.75 - 0.5 \cdot \mathbb{I}_{(\text{accepted} = 0)} + 
\begin{cases} 
25|\Delta NV| + 100 \max\left(0, \frac{\Delta TD}{TD_{before}}\right) & \text{nếu } \Delta NV < 0 \\ 
-15|\Delta NV| & \text{nếu } \Delta NV > 0 \\ 
250 \left(\frac{\Delta TD}{TD_{before}}\right) & \text{nếu } \Delta NV = 0 \text{ và } \Delta TD > 0 \\ 
0 & \text{ngược lại} 
\end{cases}
\]

(Trong đó $\mathbb{I}_{(\text{accepted} = 0)}$ là hàm chỉ thị, phạt thêm nếu số bước di chuyển bị từ chối).

\section{Kết quả - Thảo luận}
Quá trình huấn luyện trên bộ dữ liệu Solomon RC1 và RC2 chứng minh DDQN-ALNS hoạt động như một \textit{Plateau Controller} hiệu quả. Tại các điểm chững (iteration $> 300$), mạng DDQN liên tục kích hoạt các chiến lược \textit{Diversify} và \textit{Intensify}, giúp chi phí (cost) phá vỡ ngưỡng cực trị địa phương và vượt qua điểm chuẩn (BKS). Đặc biệt, mô hình thể hiện khả năng Zero-shot Transfer Learning xuất sắc khi duy trì mức độ sai số (Gap\%) chỉ 0.70\% trên cấu trúc đồ thị mới lạ hoàn toàn.

\section{Kết luận - Đề nghị}
Công trình đã chứng minh tính hiệu quả của việc kết hợp Học tăng cường sâu như một bộ não điều phối cho Metaheuristics. Hệ thống không chỉ khắc phục được hạn chế kẹt nghiệm mà còn giải quyết triệt để lỗi lạm dụng phương tiện giao hàng. Hướng nghiên cứu tiếp theo sẽ tích hợp Mạng Nơ-ron Đồ thị (GNN) để tự động trích xuất đặc trưng hình học, đồng thời mở rộng thử nghiệm trên các tập dữ liệu có quy mô lớn hơn như Gehring \& Homberger.

% ================= TÀI LIỆU THAM KHẢO =================
\newpage
\begin{thebibliography}{}

\bibitem[Dantzig và Ramser (1959)]{Dantzig1959}
Dantzig, G. B. và Ramser, J. H. (1959). \textit{The Truck Dispatching Problem}. Management Science, 6(1), 80-91.

\bibitem[Solomon (1987)]{Solomon1987}
Solomon, M. M. (1987). \textit{Algorithms for the Vehicle Routing and Scheduling Problems with Time Window Constraints}. Operations Research, 35(2), 254-265.

\end{thebibliography}

\end{document}
